% GENERATED FILE. Edit canonical lessons, then run npm run generate:books.
% canonical-source-hash: fnv1a64:69a2461b1fc031e1
% canonical-lessons: PA-C328-hausla, PA-C328-pareshan, PA-C328-tasalli, PA-C328-dabauna, PA-C328-sahina

\chapter{Caring for the Sick}
\label{ch:pa-a2-caring-for-the-sick}
\hlchaptermodality{328}

\begin{chapteropening}
\textbf{By the end of this chapter:} \emph{I can say who cared for a sick person and how.}

Complete the last lesson of chapter 328: I can say who cared for a sick person and how.
\end{chapteropening}

\section[\texorpdfstring{\pa{ਹੌਸਲਾ} (hauslā)}{ਹੌਸਲਾ (hauslā)}]{\pa{ਹੌਸਲਾ} (hauslā) — courage}

\label{lesson:PA-C328-hausla}



\begin{quote}
Before the new one: say the Punjabi for sorrow, then the Punjabi for a teardrop.
\end{quote}

\subsection*{You'll want to know: \pa{ਹੌਸਲਾ}}
\textbf{\pa{ਹੌਸਲਾ}} — \emph{hauslā} — \textquotedblleft{}courage\textquotedblright{}.

\begin{quote}
\pa{ਮਾਤਾ} \pa{ਜੀ} \pa{ਮਰੀਜ਼} \pa{ਨੂੰ} \pa{ਹੌਸਲਾ} \pa{ਦਿੰਦੇ} \pa{ਹਨ।}
\end{quote}

\begin{quote}
\emph{mātā jī marīz nū̃ hauslā dinde han.}
\end{quote}

\begin{quote}
Mother keeps the patient's spirits up.
\end{quote}

A masculine noun.

\begin{cousinweb}
\textbf{\pa{ਹੌਸਲਾ}} (\emph{hauslā}) is a loan, usually traced through Persian to Arabic \emph{ḥawṣala}, \textquotedblleft{}capacity\textquotedblright{}.
\end{cousinweb}

\begin{grammarlens}[title={What you've built}]
One more word for this chapter.
\end{grammarlens}

\subsection*{Your turn}
\begin{itemize}
\raggedright
  \item \emph{Say it:} \emph{hauslā}
  \item \emph{Say it:} \emph{hauslā}, once more
  \item \emph{Say it:} \emph{mātā jī marīz nū̃ hauslā dinde han.} — Mother keeps the patient's spirits up.
  \item \emph{Recall:} say the Punjabi for sorrow, then the Punjabi for a teardrop, then say \emph{hauslā} again
\end{itemize}

\begin{tcolorbox}[breakable,colback=teal!4,colframe=teal!35!black,title={Before you move on}]
What does \pa{ਹੌਸਲਾ} mean? (\textquotedblleft{}Courage\textquotedblright{}.) Say it once more.
\end{tcolorbox}
\section[\texorpdfstring{\pa{ਪਰੇਸ਼ਾਨ} (pareshān)}{ਪਰੇਸ਼ਾਨ (pareshān)}]{\pa{ਪਰੇਸ਼ਾਨ} (pareshān) — upset}

\label{lesson:PA-C328-pareshan}



\begin{quote}
Before the new one: say the Punjabi for a teardrop, then the Punjabi for courage.
\end{quote}

\subsection*{You'll want to know: \pa{ਪਰੇਸ਼ਾਨ}}
\textbf{\pa{ਪਰੇਸ਼ਾਨ}} — \emph{pareshān} — \textquotedblleft{}upset\textquotedblright{}.

\begin{quote}
\pa{ਪੁੱਤਰ} \pa{ਦੇ} \pa{ਬੁਖ਼ਾਰ} \pa{ਤੋਂ} \pa{ਪਿਤਾ} \pa{ਜੀ} \pa{ਪਰੇਸ਼ਾਨ} \pa{ਸਨ।}
\end{quote}

\begin{quote}
\emph{puttar de bukhār ton pitā jī pareshān san.}
\end{quote}

\begin{quote}
Father was upset about his son's fever.
\end{quote}

\begin{cousinweb}
\textbf{\pa{ਪਰੇਸ਼ਾਨ}} (\emph{pareshān}) is a loan, usually traced to Persian \emph{parēshān}, \textquotedblleft{}scattered, distressed\textquotedblright{}.
\end{cousinweb}

\begin{grammarlens}[title={What you've built}]
One more word for this chapter.
\end{grammarlens}

\subsection*{Your turn}
\begin{itemize}
\raggedright
  \item \emph{Say it:} \emph{pareshān}
  \item \emph{Say it:} \emph{pareshān}, once more
  \item \emph{Say it:} \emph{puttar de bukhār ton pitā jī pareshān san.} — Father was upset about his son's fever.
  \item \emph{Recall:} say the Punjabi for a teardrop, then the Punjabi for courage, then say \emph{pareshān} again
  \item \emph{Write it:} copy \textbf{\pa{ਪਰੇਸ਼ਾਨ}} once, checking the whole word
\end{itemize}

\begin{tcolorbox}[breakable,colback=teal!4,colframe=teal!35!black,title={Before you move on}]
What does \pa{ਪਰੇਸ਼ਾਨ} mean? (\textquotedblleft{}Upset\textquotedblright{}.) Say it once more.
\end{tcolorbox}
\section[\texorpdfstring{\pa{ਤਸੱਲੀ} (tasallī)}{ਤਸੱਲੀ (tasallī)}]{\pa{ਤਸੱਲੀ} (tasallī) — reassurance}

\label{lesson:PA-C328-tasalli}



\begin{quote}
Before the new one: say the Punjabi for courage, then the Punjabi for upset.
\end{quote}

\subsection*{You'll want to know: \pa{ਤਸੱਲੀ}}
\textbf{\pa{ਤਸੱਲੀ}} — \emph{tasallī} — \textquotedblleft{}reassurance\textquotedblright{}.

\begin{quote}
\pa{ਡਾਕਟਰ} \pa{ਦੇ} \pa{ਜਵਾਬ} \pa{ਤੋਂ} \pa{ਮੈਨੂੰ} \pa{ਤਸੱਲੀ} \pa{ਹੋਈ।}
\end{quote}

\begin{quote}
\emph{ḍākṭar de javāb ton mainū̃ tasallī hoī.}
\end{quote}

\begin{quote}
The doctor's answer reassured me.
\end{quote}

A feminine noun.

\begin{cousinweb}
\textbf{\pa{ਤਸੱਲੀ}} (\emph{tasallī}) is a loan, usually traced through Persian to Arabic \emph{tasallī}, \textquotedblleft{}consolation\textquotedblright{}.
\end{cousinweb}

\begin{grammarlens}[title={What you've built}]
One more word for this chapter.
\end{grammarlens}

\subsection*{Your turn}
\begin{itemize}
\raggedright
  \item \emph{Say it:} \emph{tasallī}
  \item \emph{Say it:} \emph{tasallī}, once more
  \item \emph{Say it:} \emph{ḍākṭar de javāb ton mainū̃ tasallī hoī.} — The doctor's answer reassured me.
  \item \emph{Recall:} say the Punjabi for courage, then the Punjabi for upset, then say \emph{tasallī} again
  \item \emph{Write it:} copy \textbf{\pa{ਤਸੱਲੀ}} once, checking the whole word
\end{itemize}

\begin{tcolorbox}[breakable,colback=teal!4,colframe=teal!35!black,title={Before you move on}]
What does \pa{ਤਸੱਲੀ} mean? (\textquotedblleft{}Reassurance\textquotedblright{}.) Say it once more.
\end{tcolorbox}
\section[\texorpdfstring{\pa{ਦਬਾਉਣਾ} (dabāuṇā)}{ਦਬਾਉਣਾ (dabāuṇā)}]{\pa{ਦਬਾਉਣਾ} (dabāuṇā) — to press gently}

\label{lesson:PA-C328-dabauna}



\begin{quote}
Before the new one: say the Punjabi for upset, then the Punjabi for reassurance.
\end{quote}

\subsection*{You'll want to know: \pa{ਦਬਾਉਣਾ}}
\textbf{\pa{ਦਬਾਉਣਾ}} — \emph{dabāuṇā} — \textquotedblleft{}to press gently\textquotedblright{}.

\begin{quote}
\pa{ਰਾਤ} \pa{ਨੂੰ} \pa{ਮੈਂ} \pa{ਦਾਦੀ} \pa{ਦੇ} \pa{ਪੈਰ} \pa{ਦਬਾਉਂਦਾ} \pa{ਹਾਂ।}
\end{quote}

\begin{quote}
\emph{rāt nū̃ maiṁ dādī de pair dabāundā hāṁ.}
\end{quote}

\begin{quote}
At night I press my grandmother's feet to ease them.
\end{quote}

Past (I …): \pa{ਮੈਂ} \pa{ਦਾਦਾ} \pa{ਜੀ} \pa{ਦਾ} \pa{ਸਿਰ} \pa{ਦਬਾਇਆ}, \emph{maiṁ dādā jī dā sir dabāiā}.

\begin{cousinweb}
\textbf{\pa{ਦਬਾਉਣਾ}} (\emph{dabāuṇā}) is inherited from a reconstructed \emph{*dabb}, \textquotedblleft{}press\textquotedblright{}; CDIAL lists Punjabi \emph{dabāuṇā}, \textquotedblleft{}to press down\textquotedblright{}, under it (CDIAL 6173).
\end{cousinweb}

\begin{grammarlens}[title={What you've built}]
One more word for this chapter.
\end{grammarlens}

\subsection*{Your turn}
\begin{itemize}
\raggedright
  \item \emph{Say it:} \emph{dabāuṇā}
  \item \emph{Say it:} \emph{dabāuṇā}, once more
  \item \emph{Say it:} \emph{rāt nū̃ maiṁ dādī de pair dabāundā hāṁ.} — At night I press my grandmother's feet to ease them.
  \item \emph{Recall:} say the Punjabi for upset, then the Punjabi for reassurance, then say \emph{dabāuṇā} again
  \item \emph{Write it:} copy \textbf{\pa{ਦਬਾਉਣਾ}} once, checking the whole word
\end{itemize}

\begin{tcolorbox}[breakable,colback=teal!4,colframe=teal!35!black,title={Before you move on}]
What does \pa{ਦਬਾਉਣਾ} mean? (\textquotedblleft{}To press gently\textquotedblright{}.) Say it once more.
\end{tcolorbox}
\section[\texorpdfstring{\pa{ਸਹਿਣਾ} (sahiṇā)}{ਸਹਿਣਾ (sahiṇā)}]{\pa{ਸਹਿਣਾ} (sahiṇā) — to put up with}

\label{lesson:PA-C328-sahina}



\begin{quote}
Before the new one: say the Punjabi for reassurance, then the Punjabi for to press gently.
\end{quote}

\subsection*{You'll want to know: \pa{ਸਹਿਣਾ}}
\textbf{\pa{ਸਹਿਣਾ}} — \emph{sahiṇā} — \textquotedblleft{}to put up with\textquotedblright{}.

\begin{quote}
\pa{ਦਾਦੀ} \pa{ਬਹੁਤ} \pa{ਦਰਦ} \pa{ਸਹਿੰਦੀ} \pa{ਹੈ} \pa{ਪਰ} \pa{ਨਹੀਂ} \pa{ਰੋਂਦੀ} \pa{ਹੈ।}
\end{quote}

\begin{quote}
\emph{dādī bahut darad sahindī hai par nahīṁ rondī hai.}
\end{quote}

\begin{quote}
Grandmother puts up with a lot of pain, but she does not cry.
\end{quote}

Past (I …): \pa{ਮੈਂ} \pa{ਬਹੁਤ} \pa{ਦਰਦ} \pa{ਸਹਿਆ}, \emph{maiṁ bahut darad sahiā}.

\begin{cousinweb}
\textbf{\pa{ਸਹਿਣਾ}} (\emph{sahiṇā}) is inherited from Sanskrit \emph{sahate}, \textquotedblleft{}conquers; bears, suffers\textquotedblright{}; CDIAL lists Punjabi \emph{sahiṇā}, \textquotedblleft{}to bear, put up with\textquotedblright{}, under it (CDIAL 13304).
\end{cousinweb}

\begin{grammarlens}[title={What you've built}]
One more word for this chapter.
\end{grammarlens}

\subsection*{Your turn}
\begin{itemize}
\raggedright
  \item \emph{Say it:} \emph{sahiṇā}
  \item \emph{Say it:} \emph{sahiṇā}, once more
  \item \emph{Say it:} \emph{dādī bahut darad sahindī hai par nahīṁ rondī hai.} — Grandmother puts up with a lot of pain, but she does not cry.
  \item \emph{Recall:} say the Punjabi for reassurance, then the Punjabi for to press gently, then say \emph{sahiṇā} again
  \item \emph{Write it:} copy \textbf{\pa{ਸਹਿਣਾ}} once, checking the whole word
\end{itemize}

\begin{tcolorbox}[breakable,colback=teal!4,colframe=teal!35!black,title={Before you move on}]
What does \pa{ਸਹਿਣਾ} mean? (\textquotedblleft{}To put up with\textquotedblright{}.) Say it once more.
\end{tcolorbox}
